\section{Bounds for the discrete circuit-growth model and for diagonal linear networks}
\label{app:discrete}

\begin{sloppypar}
This appendix states the results of two derivations: Sections A1--A4 of \repo{theory/A-discrete-model.md}, on the discrete circuit-growth model, and Sections C2--C3 of \repo{theory/C-bridge-and-certificates.md}, on diagonal linear networks and a conditional transfer to ReLU networks. Each was refereed twice, and the results are stated as revised after the second round. Numbering in parentheses is the derivation's. Proofs are sketched here; full proofs with explicit constants are in the derivations, and check scripts are in \repo{theory/checks}. ``Literature, from memory'' marks a published result whose statement the derivation could not check against the paper.
\end{sloppypar}

\subsection{Setting, counting and the Occam bound}\label{app:discrete:occam}

\paragraph{Notation.} $\mathcal X=\cube$, and $P$ is any distribution on $\mathcal X\times\{\pm1\}$, with input marginal $P_X$. The sample $S=((x_i,y_i))_{i\le n}$ is iid from $P$. We write $\langle u,v\rangle_n=\frac1n\sum_iu(x_i,y_i)v(x_i,y_i)$, $\|\cdot\|_n$ for its norm, and $\bar u=\langle u,1\rangle_n$ (in population statements, $\bar u=\E u$). The test error is $\mathrm{err}(h)=P(h(x)\ne y)$, and $\widehat{\mathrm{err}}$ is its value on $S$. Logarithms are natural. $\mathrm{kl}(q\|p)$ is the binary relative entropy, and $\mathrm{kl}^{-1}_+(q,\epsilon)=\sup\{p\ge q:\mathrm{kl}(q\|p)\le\epsilon\}\le q+\sqrt{2q\epsilon}+2\epsilon$. A test-error bound is \emph{useful} if it is at most $0.1$, and \emph{non-vacuous} if it is below $1/2$.

\paragraph{Circuits.} A \emph{library} $G$ is a finite set of gates $\phi:[-1,1]^j\to[-1,1]$ of arity $j\le k$. It is \emph{Boolean} if every gate maps $\{\pm1\}^j$ into $\{\pm1\}$. In a circuit with $T$ gates, gate $t$ applies some $\phi_t\in G$ to sources among the $d$ inputs and gates $1,\dots,t-1$, and computes $g_t$. A predictor $h=\psi(g_1,\dots,g_T)$ has a sign readout $\mathrm{sgn}(c_0+c\cdot z)$, an output-gate readout $\pm z_T$, or a real readout. $\mathcal H_T$ is the set of predictors with $T$ gates. A \emph{learner} is any rule, possibly randomized and data-dependent, that maps $S$ to circuits and readouts. The \emph{toy library} is $\{\mathrm{id},\mathrm{AND}(\pm a,\pm b),\mathrm{XOR}\}$ up to negation. It has $|G|=6$ and $k=2$, and at $d=30$ it has 2205 depth-one candidates.

\begin{lemma}[Counting (Lemmas 1.1, 1.2 in the derivation)]\label{lem:app-disc-count}
\status{Proved}\quad There are at most $e^{W_T}$ circuits with $T$ gates, where $W_T:=\sum_{t\le T}(\log|G|+k\log(d+t-1))$. On $N$ points of $\R^T$, at most $\Phi(N,T):=2\sum_{i\le T}\binom{N-1}{i}\le2N^T$ functions have the form $\mathrm{sgn}(c_0+c\cdot z)$. Hence $\log|\mathcal H_T|\le W_T+\rho_T$. Here $\rho_T=\log\Phi(2^{\min(T,d)},T)\approx T^2\log2$ (for $T\le d$) for a Boolean library with sign readout, $\rho_T=\log\Phi(2^d,T)$ for real-valued gates, and $\rho_T=\log2$ for an output-gate readout.
\end{lemma}

\begin{theorem}[Occam bound, uniform over circuit size (Theorem A1 in the derivation)]\label{thm:app-A1}
\status{Proved}\quad Fix $\delta\in(0,1)$ and let $L_T:=\log|\mathcal H_T|+\log\frac{(T+1)(T+2)}{\delta}$. With probability at least $1-\delta$ over $S$, every $T\ge0$ and every $h\in\mathcal H_T$ satisfy four bounds, where $p=\mathrm{err}(h)$ and $\hat p=\widehat{\mathrm{err}}(h)$:
\begin{enumerate}[label=(\arabic*),nosep]
\item if $\hat p\le p$, then $\mathrm{kl}(\hat p\|p)\le L_T/n$;
\item $p\le\hat p+\sqrt{2\hat pL_T/n}+2L_T/n$;
\item $p\le\hat p+\sqrt{L_T/2n}$;
\item if $\hat p=0$, then $p\le1-e^{-L_T/n}\le L_T/n$.
\end{enumerate}
A learner's output with $T(S)$ gates lies in $\mathcal H_{T(S)}$, however the gates were chosen. So the bounds hold for any learner, at every stage and every training time, with $T$ the number of distinct gate functions in use.
\end{theorem}
\begin{proof}[Proof sketch]
For a fixed $h$, the kl-Chernoff bound gives $P(\hat p\le p,\ \mathrm{kl}(\hat p\|p)>\epsilon)\le e^{-n\epsilon}$. Take a union bound over $\mathcal H_T$ with $\epsilon=L_T/n$, and then over $T$ with weights $\delta/((T+1)(T+2))$. Items (2)--(4) follow from $\mathrm{kl}(q\|p)\ge(p-q)^2/2p$, Pinsker's inequality, and $\mathrm{kl}(0\|p)=-\log(1-p)$. Nothing about the dynamics is used. This is the Occam argument~\cite{blumer1987occam}, made uniform over $T$.
\end{proof}

\paragraph{Values} \status{Numerical}. Take the toy library at $d=30$ and $\delta=0.05$.
\begin{itemize}[nosep]
\item $L_T=14.8$, $151$ and $427$ nats at $T=1$, $10$ and $20$.
\item The realizable bound (4) is useful for $T\le2,7,20,62$ and non-vacuous for $T\le11,27,102,403$, at $n=256,1024,4096,16384$.
\item The agnostic form (3) is weak: at $T=5$ and $n=4096$ its gap is $0.088$.
\item A symmetrization route in the style of Blumer, Ehrenfeucht, Haussler and Warmuth is larger at every $T\le1499$ for $n\in\{1024,4096,16384\}$.
\end{itemize}

\paragraph{A norm version} (Theorem A1$'$, \status{Proved given three textbook Rademacher lemmas}). This covers $f=c_0+\sum_jc_jg_j$ with $g_j$ in a finite class $\mathcal F$, $\sum_j|c_j|\le B$ and $|c_0|\le1+B$. Uniformly over $B=2^j$, with probability $1-\delta$, two gaps are bounded. The clipped square loss exceeds its empirical value by at most $4\mathfrak R_n(B)+2\sqrt{\log(2/\pi_j\delta)/2n}$. The 0-1 error exceeds the empirical ramp loss at a margin $\theta$ fixed in advance by at most $\frac2\theta\mathfrak R_n(B)+\sqrt{\log(2/\pi_j\delta)/2n}$. Here $\mathfrak R_n(B)=(1+B)/\sqrt n+B\sqrt{2\log(2|\mathcal F|)/n}$ and $\pi_j=1/((j+1)(j+2))$.

\subsection{The clock race and what training time controls}\label{app:discrete:time}

\begin{definition}[Clock race (Definition 2.1 in the derivation)]\label{def:app-race}
Fix $\Theta>0$. The state at time $t$ is a circuit of \emph{fired} gates, a readout $c(t)$, and the residual $r=y-f_t$, where $f_t=c_0+\sum c_mg_m$.
\begin{itemize}[nosep]
\item A \emph{candidate} $m$ is a gate on inputs or fired gates. Its \emph{availability time} $a_m$ is the last firing time among its gate sources, or $0$ if it has none.
\item Its \emph{clock} is $S_m(t)=\int_{a_m}^tU_m$, with utility $U_m=\langle r,g_m\rangle_n$. It fires when $|S_m|=\Theta$, with sign $\sigma_m=\mathrm{sgn}\,S_m$. Several candidates may fire at once, and clocks are never reset.
\item After each event the readout is refit, in one of two variants. (LS): least squares with an intercept over all fired gates. (SC), for a depth-one dictionary only: least squares with an intercept over $\mathrm{Sat}=\{m:|S_m|=\Theta\}$, subject to $\sigma_mc_m\ge0$.
\end{itemize}
\emph{Effective time} is $\tau(t)=\int_0^t\|r\|_n=\int_0^t\sqrt{2\hat L}$, and $R(t)=\int_0^tr$. The \emph{population race} uses $L^2(P)$ in place of the sample.
\end{definition}

\paragraph{Basic facts} (Lemma 2.2, \status{Proved}).
\begin{itemize}[nosep]
\item $|U_m|\le\|r\|_n\le1$, so $\tau\le t$.
\item In (LS), a fired gate keeps $U_m=0$.
\item In (SC), $|S_m|\le\Theta$.
\item With one threshold, rescaling $\Theta$ rescales time.
\end{itemize}
(SC) results hold up to any accumulation point of event times. Whether such points exist is open.

\paragraph{Scope.}
\begin{itemize}[nosep]
\item For moves needing $k=2$ fresh weight layers, $\Theta=c\log(1/\alpha)$, with a constant $c$ that depends on the parametrization ($c=\frac12$ for the diagonal network of \cref{app:discrete:diag}).
\item Fresh chains with $k\ge3$ cost about $\alpha^{-(k-2)}$, so the first move of a three-layer network started at scale $\alpha$ costs about $\alpha^{-1}$ \status{Heuristic}. The depth-one race therefore models diagonal or 2-homogeneous networks.
\item The diagonal network's small-$\alpha$ limit~\cite{pesme2023saddle} is recalled to match (SC) up to the intercept. AGF~\cite{kunin2025agf} also accumulates utility without resets. Both are Literature, from memory.
\item (LS) is not the depth-one limit.
\item Deeper moves are \status{Heuristic}.
\end{itemize}

\begin{proposition}[Depth (Proposition A2.1 in the derivation)]\label{prop:app-A21}
\status{Proved}\quad In either variant, a gate of depth $D$ fires only when $\tau(t)\ge D\Theta$. So the depth at time $t$ is at most $\lfloor\tau(t)/\Theta\rfloor\le\lfloor t/\Theta\rfloor$. The proof is by induction on $D$, since $\Theta=|S_m(t_m)|\le\tau(t_m)-\tau(a_m)$.
\end{proposition}

\begin{proposition}[No uniform gate-count bound (Proposition A2.2 in the derivation)]\label{prop:app-A22}
\status{Proved}\quad Consider the population race with uniform inputs and $y=\mathrm{MAJ}(x_1,\dots,x_m)$, for odd $m\le d$. Let $u_m=\binom{m-1}{(m-1)/2}2^{1-m}\approx\sqrt{2/\pi m}$.
\begin{enumerate}[label=(\arabic*),nosep]
\item For $G=\{\mathrm{id}\}$, $m$ gates fire together at $\Theta/u_m$ and nothing fires earlier, so $T\approx\frac2\pi(t/\Theta)^2$.
\item For the toy library, exactly $m^2$ functions, distinct up to sign, fire together at $\Theta/u_m$: the identities, $\mathrm{AND}(x_i,x_j)$ and $\mathrm{AND}(\neg x_i,\neg x_j)$. So $T\approx(\frac2\pi)^2(t/\Theta)^4$.
\item Duplicates with equal availability times fire together, so no function of $t/\Theta$ bounds the syntactic gate count uniformly over libraries.
\end{enumerate}
\end{proposition}
\begin{proof}[Proof sketch]
$\E[yx_i]=u_m$, the chance that the other votes tie, and $\E[yx_ix_j]=0$. Expanding the AND gates as $\mathrm{AND}(\epsilon_1a,\epsilon_2b)=(\epsilon_1a+\epsilon_2b+\epsilon_1\epsilon_2ab-1)/2$ shows that the listed gates have utility $\pm u_m$, and all others at most $u_m/2$ in absolute value.
\end{proof}
So no bound $T\le C(t/\Theta)^2$ with $C$ independent of $d$ holds for the toy library uniformly over the targets $\mathrm{MAJ}_m$, $m\le d$. At $n=4096$, maj3 fires 9 gates at $t/\Theta\in[1.98,2.12]$ \status{Numerical}.

\begin{theorem}[Gate count times incoherence (Theorem A2.3 in the derivation)]\label{thm:app-A23}
\status{Proved}\quad Let $K_B$ be the centred empirical Gram matrix of a set $B$ of gates, with entries $\langle g_m-\bar g_m,g_{m'}-\bar g_{m'}\rangle_n$, and $\kappa_B=\lambda_{\max}(K_B)$.
\begin{enumerate}[label=(\arabic*),nosep]
\item In (LS), at any depth, every set $B$ fired by time $t$ satisfies $|B|\le\kappa_B(\tau(t)/\Theta)^2$. If all $a_m=0$, $\tau$ may be replaced by $\|R(t)\|_n$.
\item In (SC), $|\mathrm{Sat}(t)|\le\kappa_{\mathrm{Sat}(t)}(\|R(t)\|_n/\Theta)^2$, and $\mathrm{supp}\,c\subseteq\mathrm{Sat}$.
\end{enumerate}
So the number $T(t)$ of distinct fired functions is at most $\kappa(t)(\tau/\Theta)^2$. Let $\mu_B$ be the largest absolute centred covariance within $B$, and $x=(\tau/\Theta)^2$. If $\mu_Bx<1$, then $|B|\le(1-\mu_B)x/(1-\mu_Bx)$, by Gershgorin's theorem.
\end{theorem}
\begin{proof}[Proof sketch]
Each $m\in B$ has $\sigma_mS_m(t)=\Theta$. The refit intercept gives $\bar r=0$, so $U_m=\langle r,\tilde g_m\rangle_n$ with $\tilde g_m=g_m-\bar g_m$. With $B(s)=\{m\in B:a_m\le s\}$, $|B|\Theta=\int_0^t\langle r,\sum_{B(s)}\sigma_m\tilde g_m\rangle_n\le\tau(t)\sqrt{\kappa_B|B|}$. The last step uses the fact that a principal submatrix has no larger top eigenvalue. When all $a_m=0$, apply Cauchy--Schwarz to $\langle R,\cdot\rangle_n$ instead.
\end{proof}
The population analogue is tight up to a factor $\pi/2$ (\cref{prop:app-A22}(1)). A library-wide coherence does not help, since the toy library has $\mu=\frac12$.

\begin{theorem}[$\ell_1$ norm of the readout (Theorem A2.4 in the derivation)]\label{thm:app-A24}
\status{Proved}\quad In (SC) on a depth-one dictionary, $\Theta\|c(t)\|_1=\langle R(t),f_t\rangle_n\le\|R(t)\|_n\sqrt{1-\bar y^2}\le\tau(t)$, and $|c_0|\le1+\|c\|_1$.
\end{theorem}
\begin{proof}[Proof sketch]
KKT gives $c_mS_m=|c_m|\Theta$, with $S_m=\langle R,g_m\rangle_n$. Since $\langle R,1\rangle_n=0$, $\Theta\|c\|_1=\langle R,f_t-\bar y\rangle_n$. Here $f_t-\bar y$ is the projection of $y-\bar y$ onto a cone, so its norm is at most $\sqrt{1-\bar y^2}$. This is LASSO duality, with $R$ as the residual and $\Theta$ as the penalty.
\end{proof}

\paragraph{Checks} \status{Numerical}. These are toy races at $d=30$ with $n=64$--$4096$.
\begin{itemize}[nosep]
\item \Cref{thm:app-A1,prop:app-A21,thm:app-A23} held at every event.
\item The ratio $\kappa\tau^2/(\Theta^2T)$ is $1.00$ on par2, where the bound is attained. It is at most 7 on maj3, and3 and dnf3, and up to 400 on par3, where one noise gate fires late.
\item \Cref{thm:app-A24} holds to machine precision in the toy (SC) runs, and to $3\cdot10^{-14}$ on 40 random correlated dictionaries with 130 exits from the saturated set.
\end{itemize}

\begin{corollary}[Test error along training (Corollary A2.5 in the derivation)]\label{cor:app-A25}
\status{Proved for (a), (a$'$); proved given the lemmas of A1$'$ for (b)}\quad Let $h_t=\mathrm{sgn}f_t$ and $D(t)=\lfloor\tau(t)/\Theta\rfloor$. Each of (a), (a$'$) and (b) holds with probability $1-\delta$, simultaneously for all $t$.
\begin{enumerate}[label=(\alph*),nosep]
\item For (LS), \cref{thm:app-A1} holds for $h_t$ with $T(t)\le\lfloor\kappa(t)\tau^2/\Theta^2\rfloor$. For example, at zero training error, $\mathrm{err}(h_t)\le L_{\lfloor\kappa\tau^2/\Theta^2\rfloor}/n$.
\item[(a$'$)] For (LS) or (SC), in a race that merges candidates equal up to sign on the sample: let $\mathcal F_{\le D}$ be the non-constant depth-$\le D$ functions that are distinct up to sign on the sample, and $\bar\kappa_D$ their top centred Gram eigenvalue. These depend only on the library and the unlabelled inputs. Then \cref{thm:app-A1} holds with $T(t)\le\min\{|\mathcal F_{\le D(t)}|,\lfloor\bar\kappa_{D(t)}\tau^2/\Theta^2\rfloor\}$, and at depth one $\tau$ may be replaced by $\|R\|_n$.
\item[(b)] For (SC) at depth one, Theorem A1$'$ holds with $B\le2\max(1,\tau/\Theta)$. Its gap is therefore of order $(1+\tau/\Theta)\sqrt{\log(|G|d^k)/n}$, whatever the gate count.
\end{enumerate}
\end{corollary}
For (a$'$), the fired representatives lie in $\mathcal F_{\le D(t)}$ (\cref{prop:app-A21}), and up to sign flips their Gram matrix is a principal submatrix of the library's.

\paragraph{How strong these are.}
\begin{itemize}[nosep]
\item \emph{(a) is a decomposition, not a time-only bound.} $\kappa(t)$ belongs to the learned circuit and can grow (3 for maj3 after its first event). It is never tighter than \cref{thm:app-A1} at the realized size.
\item \emph{(a$'$) is time-only but weak.} For the toy library, $\bar\kappa_1=d$ in the population, and $35$ and $32$ empirically at $n=4096$ and $16384$. At zero training error, (a$'$) gives $0.21$ and $0.053$ at $\|R\|_n=\Theta$, and $0.61$ and $0.195$ just below $2\Theta$. \Cref{thm:app-A1} at par2's realized size gives $0.004$ and $0.001$.
\item \emph{(a$'$) needs the merging.} A referee counterexample shows that it fails without it.
\item \emph{(b) is near-vacuous.} Its 0-1 bound at $n=4096$ is $0.40$--$0.97$.
\end{itemize}
So training time bounds the depth, the incoherence-weighted size and the depth-one readout norm. It bounds the raw gate count only through $\kappa(t)$ or $\bar\kappa\approx d$.

\subsection{Minibatches: reproducing the population run}\label{app:discrete:sgd}

\paragraph{The population greedy run (PG).} PG adds one gate per step. Given $C_t$, let $f_t$ be the $L^2(P)$ projection of $y$ onto $V_t=\mathrm{span}\{1,C_t\}$, and $r_t=y-f_t$. PG adds $m_t=\arg\max_m|U_t(m)|$, where $U_t(m)=\E[r_tg_m]$. The maximum is over $M_t\le|G|(d+t)^k$ candidates, identified up to sign. Three assumptions are made:
\begin{itemize}[nosep]
\item a gap $\gamma>0$: $|U_t(m_t)|\ge|U_t(m)|+\gamma$ for $m\ne\pm m_t$ and $t<T$;
\item $|r_t|\le R$;
\item $\lambda_*=\min_{t\le T}\lambda_{\min}(K_t)>0$, for the $L^2(P)$ Gram matrices $K_t$ of $(1,C_t)$.
\end{itemize}

\paragraph{The empirical run (EG).} At step $t$, EG uses a batch $B_t$ of size $b$: either $S$ itself (full batch) or a fresh batch. It refits on $B_t$, adds a maximizer of $|\langle\hat r_t,g_m\rangle_{B_t}|$, and outputs the sign of its final refit.

\begin{theorem}[Decision stability (Theorem A3.1 in the derivation)]\label{thm:app-A31}
\status{Proved}\quad Let $\epsilon=\gamma/(2+2\sqrt{2T/\lambda_*})$, and suppose that $b\ge\frac{2R^2}{\epsilon^2}\log\frac{4\sum_{t<T}(M_t+t+1)}{\delta}$ and $b\ge\frac{8(T+1)^2}{\lambda_*^2}\log\frac{4(T+1)^3}{\delta}$. Then, for both protocols, with probability $1-\delta$:
\begin{enumerate}[label=(\arabic*),nosep]
\item EG's circuit computes $C_T$;
\item if $y\in V_T$ almost surely, EG's predictor equals $y$ almost surely.
\end{enumerate}
With probability $1-2\delta$, \cref{thm:app-A1} holds for EG's predictor, on the sample of its final refit, with $L=\rho_T+\log(1/\delta)$ and no wiring term.
\end{theorem}
\begin{proof}[Proof sketch]
The functions $r_t$, $g_m$ and $K_t$ belong to the population run, so Hoeffding's inequality and a union bound give $\epsilon$-accurate empirical utilities, and $\lambda_{\min}(\hat K_t)\ge\lambda_*/2$. Inductively, if EG holds $C_t$, then $\hat r_t=(I-\hat\Pi_t)r_t$ with $\|\hat\Pi_tr_t\|^2<2(t+1)\epsilon^2/\lambda_*$. So every empirical utility is within $\gamma/2$ of $U_t$, and the gap fixes the argmax.
\end{proof}

\paragraph{Related results.}
\begin{itemize}[nosep]
\item \emph{Reuse} (Theorem A3.5, \status{Proved}). If batches are drawn with replacement from a fixed $S$ of size $n$, the same conclusions hold, with the Occam bound on $S$. The conditions use $(n^{-1/2}+b^{-1/2})^{-2}$ in place of $b$, and $\delta/2$ in place of $\delta$.
\item \emph{Thresholds} (\status{Numerical}). For par2, with $\mathrm{XOR}\in G$, $\gamma=\frac12$, $R=\lambda_*=1$ and $M_0=2205$, the theorem needs $b\ge2254$. With reuse at $n=16384$ it needs $b\ge6225$. So the experiments' $b=64$ is not covered. At $T=30$ and $\lambda_*=\frac12$, the Gram condition needs $4.5\cdot10^5$ samples.
\item \emph{Limits.} Stability gives exact recovery with no factor $1/\varepsilon$ for a target error $\varepsilon$, but it pays $1/\gamma^2$ and says nothing under ties or zero utilities.
\item \emph{The SGD race.} For the race, only clock concentration is proved (Lemma A3.2, \status{Proved}). Under fresh batches, with residuals bounded by $R$ and a fixed finite candidate set $\mathcal M$, with probability $1-\delta$ every clock stays within $R\sqrt{2\eta t_0\log(2|\mathcal M|/\delta)/b}$ of its drift clock up to time $t_0$; the observed errors are 75--88\% of this \status{Numerical}. That the race reproduces the population race is \status{Heuristic}.
\item \emph{Online-to-batch} (Proposition A3.3, \status{Proved}; cf.~\cite{cesabianchi2004online}). In one-pass training on fresh samples, pre-update minibatch losses estimate the trajectory's average test loss.
\end{itemize}

\subsection{Greedy growth versus minimum circuits}\label{app:discrete:greedy}

\begin{proposition}[Equal, overshoot, never (Propositions A4.1--A4.3 in the derivation)]\label{prop:app-A4}
\status{Proved; item 2 for the population race only}\quad Let $s^*$ be the size of the smallest consistent circuit.
\begin{enumerate}[label=(\arabic*),nosep]
\item If $y=\pm g^*$ for a depth-one candidate $g^*$, and $0<P(y=1)<1$, then $g^*$ and its duplicates fire first, at $\Theta/\mathrm{Var}(y)$. The race then stops with $T=1=s^*$. The same holds on a sample containing both labels, with $\mathrm{Var}_n$, and with duplicates meaning candidates equal to $\pm y$ on the sample.
\item Take uniform inputs, the toy library, and $y=\mathrm{maj}_3$. Then 9 gates fire at $2\Theta$, and the 0-1 error is already zero. Exactly 12 depth-two gates of utility $\frac14$ fire next, at $6\Theta$. So the race ends with 21 gates, while $s^*=3$ (\status{Numerical}, by exhaustive search). For $\mathrm{MAJ}_m$, $m^2$ gates fire where $m$ suffice.
\item Take uniform inputs and $y=\chi_A$ with $|A|=s>k$. Then nothing ever fires in the population race. On a sample, with probability $1-\delta$, nothing fires before $\Theta\sqrt{n/(8\log(2(M_1+1)/\delta))}$, where $M_1$ is the number of depth-one candidates (Proposition A2.6). If $\mathrm{XOR}\in G$, then with probability $1-\delta$ the minimum-consistent-circuit learner with an output-gate readout has error at most $L_{s-1}/n=O((s\log(d+s)+\log\frac1\delta)/n)$.
\end{enumerate}
\end{proposition}
\begin{proof}[Proof sketch]
(1) For $|g|\le1$, $|\E[(y-\bar y)g]|\le\E|y-\bar y|=\mathrm{Var}(y)$, with equality only for $g=\pm y$. (2) The 9 gates span the degree-$\le2$ functions, which leaves the residual $-\frac12x_1x_2x_3$, and a Fourier lemma bounds the depth-two utilities by $\frac14$. (3) Depth-one gates see fewer than $s$ variables, and the error bound is \cref{thm:app-A1}(4).
\end{proof}
So both learners get the same Occam bound as a function of $T$, and they differ in which $T$ they reach \status{Numerical}.
\begin{itemize}[nosep]
\item maj3 reaches zero training error with 6 gates at $n\ge256$.
\item dnf3 needs 18 gates where 6 suffice. Its bound at $n=4096$ is $0.084$, where 6 gates would give $0.019$.
\item par3 fires only noise gates, at $t/\Theta=2.65$--$20.0$ (the proved lower bound is $0.84$--$6.71$). Its bound of $0.52$--$0.95$ certifies nothing.
\end{itemize}

\subsection{Diagonal linear networks (Section C2 of the derivation)}\label{app:discrete:diag}

\paragraph{Setting.} The dictionary consists of $D$ characters $\chi_J(x)=\prod_{i\in J}x_i$; for example, taking all $|J|\le k$ gives $D=D_k=\sum_{i\le k}\binom di$. Let $X_{iJ}=\chi_J(x_i)$ and $K=X^\top X/n$. The model is $f_w=\sum_Jw_J\chi_J$, with $w=u\odot u-v\odot v$ and loss $L(w)=\frac1{2n}\|Xw-y\|^2$. It is trained by gradient flow from $u=v=\alpha\mathbf1$. Write $r=y-Xw$, $g=\nabla L=-X^\top r/n$, $G=\int_0^tg$, $\tau=\int_0^t\|r\|_n$ and $R=\int_0^tr$, and set $\Theta_\epsilon=\frac14\operatorname{arsinh}\frac\epsilon{2\alpha^2}\ge\frac14\log\frac\epsilon{\alpha^2}$.

\paragraph{Closed form and no early growth} (Lemma C2.1, Theorem C2.2, \status{Proved}).
\begin{itemize}[nosep]
\item $w(t)=-2\alpha^2\sinh(4G(t))$. So $-G_J$ is exactly the race clock, and $|w_J|\ge\epsilon$ exactly when it reaches $\Theta_\epsilon\approx\frac12\log\frac1\alpha$.
\item Since $|G_J|\le\tau$, nothing reaches $\epsilon$ while $\tau(t)<\Theta_\epsilon$.
\item Let $U_0=\max_J|\langle y,\chi_J\rangle_n|$ and $\epsilon_D=\sqrt{\lambda_{\max}(K)D}\,\epsilon$. If $U_0>\epsilon_D$, the first escape time lies in $[\Theta_\epsilon/(U_0+\epsilon_D),\Theta_\epsilon/(U_0-\epsilon_D)]$.
\end{itemize}

\begin{theorem}[How many coordinates can have grown (Theorem C2.3 in the derivation)]\label{thm:app-C23}
\status{Proved}\quad Let $B=\{J:|w_J(t)|\ge\epsilon\}$, $N_\epsilon=|B|$ and $\kappa_B=\lambda_{\max}(X_B^\top X_B/n)$.
\begin{enumerate}[label=(\arabic*),nosep]
\item $N_\epsilon\le(\int_0^t\|\nabla L\|_2)^2/\Theta_\epsilon^2$.
\item $N_\epsilon\le\kappa_B\tau^2/\Theta_\epsilon^2$. Let $\mu_n=\max_{J\ne J'}|\langle\chi_J,\chi_{J'}\rangle_n|$ and $x=(\tau/\Theta_\epsilon)^2$. Then $N_\epsilon\le(1-\mu_n)x/(1-\mu_nx)$ when $\mu_nx<1$. For uniform inputs and characters of degree at most $k$, with probability $1-\delta'$, $\mu_n\le\bar\mu:=\sqrt{2\log(2D_{\min(2k,d)}/\delta')/n}$.
\item $\sum_J|w_J||G_J|=\langle R,f_w\rangle_n$, so $\sum_{J\in B}|w_J|\le2\tau/\Theta_\epsilon$.
\end{enumerate}
\end{theorem}
\begin{proof}[Proof sketch]
Each $J\in B$ has $|G_J|\ge\Theta_\epsilon$, and (1) follows from Minkowski's inequality. For (2), with $\sigma_J=-\mathrm{sgn}\,G_J$, $N\Theta_\epsilon\le\int_0^t\langle r,\sum_B\sigma_J\chi_J\rangle_n\le\tau\sqrt{\kappa_BN}$, as in \cref{thm:app-A23}. The bound on $\mu_n$ is Hoeffding's inequality for the characters $\chi_{J\triangle J'}$.
\end{proof}
Item (2) is exact at the first crossing for a single-character target that is empirically orthogonal to the rest of the dictionary. At nonzero coherence it is strict; the derivation sketches, without a full proof, that it becomes tight as $\alpha\to0$.

\begin{theorem}[Generalization uniform in $t$ and $\epsilon$ (Theorem C2.4 in the derivation)]\label{thm:app-C24}
\status{Proved}\quad Let $\ell_N=\log\binom DN+\log\Phi(2^{\min(N,d+1)},N)$ and $\mu_\epsilon(t)=\sum_{J\notin B}|w_J(t)|$. With probability $1-\delta$, for all $t\ge0$ and $\epsilon>0$, with $N=N_\epsilon(t)$:
\[P(yf_{w(t)}\le0)\le\mathrm{kl}^{-1}_+\big(\widehat P_S[yf_{w(t)}\le2\mu_\epsilon(t)],\ (\ell_N+\log\tfrac{(N+1)(N+2)}\delta)/n\big).\]
\end{theorem}
\begin{proof}[Proof sketch]
Split $f_w=f_L+f_s$, with $f_L$ supported on $B$ and $\|f_s\|_\infty\le\mu$. The set $\{yf_L\le\mu\}$ is a threshold of $(y\chi_J(x))_{J\in B}$, which takes at most $2^{\min(N,d+1)}$ values. By Cover's count there are therefore at most $e^{\ell_N}$ such sets. Apply kl-Chernoff with a union bound over the sets and over $N$, and compare margins on test and training points.
\end{proof}

\begin{corollary}[Time-indexed form (Corollary C2.5 in the derivation)]\label{cor:app-C25}
\status{Proved}\quad Fix $\gamma\in(0,1]$. Let $\epsilon=\gamma/2D$, $\Theta_\gamma=\frac14\operatorname{arsinh}\frac{\gamma}{4D\alpha^2}$, $x=(\tau/\Theta_\gamma)^2$ and $\ell^*_N=\max_{N'\le N}\ell_{N'}$. On the event of \cref{thm:app-C24}, every $\bar N\in[N_\epsilon(t),D]$, possibly data-dependent, satisfies
\[\mathrm{err}(\mathrm{sgn}f_{w(t)})\le\mathrm{kl}^{-1}_+\big(\widehat P_S[yf_{w(t)}\le\gamma],\ (\ell^*_{\bar N}+\log\tfrac{(\bar N+1)(\bar N+2)}\delta)/n\big).\]
There are three admissible choices of $\bar N$, each capped at $D$:
\begin{enumerate}[label=(\arabic*),nosep]
\item $\lfloor\kappa_Bx\rfloor$, with $\kappa_B$ measured;
\item $\lfloor(1-\mu_n)x/(1-\mu_nx)\rfloor$, if $\mu_nx<1$;
\item for uniform inputs and all characters of degree $\le k$, the same with $\bar\mu$ in place of $\mu_n$, and optionally $t$ in place of $\tau$, if $\bar\mu x<1$, at total probability $1-\delta-\delta'$.
\end{enumerate}
\end{corollary}
Only choice (3) depends on $(n,D,\tau)$ alone, and it needs $n>2\log(2D_{\min(2k,d)}/\delta')(\tau/\Theta_\gamma)^4$. $D$ enters only through $\log D$.

\paragraph{Discrete GD and SGD} (Proposition C2.6, \status{Proved}). Suppose $2\eta\bar G<1$, where $\bar G=\max_j\|g_j\|_\infty$.
\begin{itemize}[nosep]
\item No coordinate grows before the minibatch effective time $\eta\sum_j\|r_j\|_{B_j}$ reaches $\Theta_\epsilon(1-4\eta^2\bar G^2)$.
\item Under SGD, the count is bounded only through the integrated gradient norms $\eta\sum_j\|g_j\|_2$, which is loose.
\item Under full-batch GD, the $\kappa_B\tau^2$ form holds, with $\Theta_\epsilon$ reduced by an explicit $O(\eta^2\bar G^2\tau)$.
\end{itemize}

\paragraph{Checks} \status{Numerical}. These are at $d=10$, $k=3$, $D=176$, with exact test errors. The generalization bounds are at $t=60$ and $\gamma=\frac12$.
\begin{itemize}[nosep]
\item \emph{Size bounds.} The worst ratio of $N$ to $\kappa_B\tau^2/\Theta^2$ is $0.33$--$0.94$, and $1.00$ for par2 at $\alpha=10^{-8}$. The discrete counts reach $0.62$--$0.80$ of their bound under GD, and only $0.07$--$0.12$ under SGD.
\item \emph{A posteriori.} At $\alpha=10^{-3}$, \cref{thm:app-C24} gives $0.107$ for par2 ($n=100$), $0.275$ and $0.077$ for maj3 ($n=100$ and $400$), and $0.166$ for and3 ($n=400$).
\item \emph{Time-indexed.} Choice (1) gives $0.218$, $0.810$, $0.316$ and $0.477$ on the same runs. Choices (2) and (3) are vacuous except for par2 at $\alpha=10^{-8}$, where they give $0.107$ and $0.163$. So the time-only form works here only for a single character at tiny $\alpha$.
\item \emph{Off-dictionary targets.} maj5 lies outside the dictionary and gets $0.951$, against a true error of $0.073$. par4 is uncorrelated with the dictionary and gets $0.988$.
\end{itemize}

\subsection{Conditional transfer to ReLU networks (Section C3 of the derivation)}\label{app:discrete:relu}

\paragraph{Known results and growth lemmas.} Rigorous results cover early alignment and the first escape in two-layer and homogeneous networks, for example~\cite{boursier2022early}. Exact saddle-to-saddle limits are known only for diagonal networks~\cite{pesme2023saddle,kunin2025agf} (Literature). The derivation proves one-sided growth bounds (Lemmas C3.1, C3.2, \status{Proved}). Let $\lambda_X$ be the top eigenvalue of the second moment of $(x,1)$.
\begin{itemize}[nosep]
\item \emph{Two layers.} In a two-layer ReLU network under gradient flow or descent, each unit's $\rho_j=a_j^2+\|w_j\|^2$ satisfies $\rho_j(t)\le\rho_j(0)e^{2\sqrt{\lambda_X}\tau(t)}$, with $\tau$ the discrete effective time $\eta\sum_k\|r_k\|_n$ for descent.
\item \emph{Deep networks.} In a bias-free deep ReLU network with $L\ge3$ under gradient flow, all layers grow equally~\cite{du2018algorithmic}. The quantity $Q=\max_l\|W_l(0)\|_F^2+\Delta$ satisfies $Q(t)^{1-L/2}\ge Q(0)^{1-L/2}-(L-2)\sqrt{\lambda_X}\tau$. So, with layers of Frobenius norm $\asymp\alpha$, escape takes effective time at least of order $\alpha^{-(L-2)}$.
\item \emph{Tightness} (\status{Numerical}). At $L=3$ this is loose by $14$--$21\times$, and at $L=4$ it is vacuous.
\item \emph{Scope.} The study's $\mu$P Adam networks are outside both lemmas.
\end{itemize}

\paragraph{The hypothesis.} Let $f_{\theta(t)}$ be the network at checkpoints $t\in\mathcal T$. Run the (LS) race on the same $S$ with a finite library. Let $h^M_s=\mathrm{sgn}f^M_s$ be its state at race time $s$, with $T^M(s)$ gates, and let $s(t)$ be a clock calibration. The \emph{saddle-to-saddle approximation hypothesis} (SSA) has two forms:
\begin{itemize}[nosep]
\item sign form: $P_X(\mathrm{sgn}f_{\theta(t)}\ne h^M_{s(t)})\le\eta(t)$;
\item $L^2$ form: $\|\mathrm{clip}f_{\theta(t)}-\mathrm{clip}f^M_{s(t)}\|_{L^2(P_X)}\le\eta(t)$.
\end{itemize}
For a Boolean library, let $L^{\rm sgn}_T=W_T+\log\Phi(2^{\min(T,d)},T)$ and $L^{\rm mar}_T=W_T+\log\Phi(2^{\min(T+1,d+1)},T+1)$. For real-valued gates the exponents are $d$ and $d+1$.

\begin{theorem}[The network inherits the race bound (Theorem C3.4 in the derivation)]\label{thm:app-C34}
\status{Conditional: proved under SSA}\quad There is an event of probability $1-\delta$ over $S$, not depending on the network, on which the following hold for every $t\in\mathcal T$. Here $s=s(t)$ and $T=T^M(s)\le\kappa^M(s)(\tau^M(s)/\Theta)^2$.
\begin{enumerate}[label=(\arabic*),nosep]
\item Under the sign form, $\mathrm{err}(\mathrm{sgn}f_{\theta(t)})\le\eta(t)+\mathrm{kl}^{-1}_+\big(\widehat{\mathrm{err}}_S(h^M_s),(L^{\rm sgn}_T+\log\frac{2(T+1)(T+2)}\delta)/n\big)$.
\item Under the $L^2$ form, for every $\gamma\in(0,1)$, $\mathrm{err}(\mathrm{sgn}f_{\theta(t)})\le\frac{\eta(t)^2}{\gamma^2}+\mathrm{kl}^{-1}_+\big(\widehat P_S[yf^M_s\le\gamma],(L^{\rm mar}_T+\log\frac{2(T+1)(T+2)}\delta)/n\big)$.
\end{enumerate}
\end{theorem}
\begin{proof}[Proof sketch]
Before training, fix two Occam events, each at level $\delta/2$: one for threshold-readout circuits, and one for margin sets of circuits. The sign form then adds the disagreement. In the $L^2$ form, $yf_\theta\le0$ forces either $yf^M\le\gamma$ or a clipped gap larger than $\gamma$, and Chebyshev's inequality bounds the latter.
\end{proof}

\paragraph{What this gives.}
\begin{itemize}[nosep]
\item \emph{A conversion only.} The theorem converts the hypothesis into a bound, and its size term is data-dependent through $\kappa^M$.
\item \emph{Testing SSA} (Corollary C3.5, \status{Proved}). Held-out unlabelled inputs make the sign form checkable. The distillation certificate (\cref{thm:C1}) then bounds the network with no hypothesis at all. So SSA is needed only for the predictive statement.
\item \emph{A first test} (\status{Numerical}). On the saved $\mu$P Adam networks, and3 and dnf3 match the race's constant state on their first plateau, which is trivial. In transitions, the best disagreement with any race state is $0.1$--$0.3$ (dnf3, mux, stair5), or no better than the disagreement with the target (maj3, majpar, parmaj). On par3 the race never starts, while the network learns.
\end{itemize}
So \cref{thm:app-C34} cannot yet be applied to the study's runs, and proving SSA beyond the first escape is open.
